\section{Finite measurement precision and the exact crossover}
\label{sec:finite-precision}
Suppose the designer observes the offset bin containing $c$, rather
than its exact value. For an integer $N\ge1$, let $\Delta=4/N$ and
partition $[-2,2]$ at $-2+j\Delta$, $j=0,\ldots,N$. The bins are
half-open, except at the final endpoint. Write $\mathcal I_\Delta$
for this partition and define
\begin{equation}\label{eq:bin-value}
 \mathcal Q(M,\Delta)=
 \inf_{\substack{D_I\in L^1(\Gamma),\ D_I\ge0\\
                  \int_\Gamma D_I=M\ \text{for every }I}}
       \frac14\sum_{I\in\mathcal I_\Delta}\int_I J_c(D_I)\dd c.
\end{equation}
The budget is available separately for every observed bin; it is
not shared across bins. There are finitely many decisions at each
$M,\Delta$, so this is a measurable policy in the sense of
\cref{sec:information-transfer}. Its value is the sum of the
conditional infima, multiplied by $\Delta/4$: choose an arbitrarily
close competitor in each of the finitely many bins to obtain the
reverse inequality. This does not require existence of a circle
optimizer.

The next theorem includes the intermediate regime, in which the
measurement error and the optimal placement width have comparable
size. Its coefficient is expressed through a fully specified local
optimization problem, rather than an elementary formula.

\begin{theorem}[Measurement precision and mobility budget]
\label{thm:precision}
For the cosine ensemble, jointly as $M\downarrow0$ and
$\Delta=4/N\downarrow0$,
\begin{equation}\label{eq:precision-order}
 \mathcal Q(M,\Delta)\asymp
       M^{-1/5}+M^{-1/4}\Delta^{1/4},
\end{equation}
with constants independent of the relative rates of the two limits.
Let $\mathcal F_{\rm ctr}$ be the local value in
\eqref{eq:center-value} below and put
\begin{equation}\label{eq:crossover-function}
 \mathcal H(\tau)=\frac{2^{6/5}}4\int_{-1}^1
  (1-c^2)^{-4/5}
  \mathcal F_{\rm ctr}\!\left(2^{1/5}\tau(1-c^2)^{-3/10}\right)\dd c,
       \qquad \tau\ge0.
\end{equation}
This integral is finite and defines a continuous function. For
every finite $\tau\ge0$,
\begin{equation}\label{eq:precision-crossover}
 \frac{\Delta}{M^{1/5}}\longrightarrow\tau
 \quad\Longrightarrow\quad
 M^{1/5}\mathcal Q(M,\Delta)\longrightarrow\mathcal H(\tau).
\end{equation}
In particular, $\mathcal H(0)=K_{\rm obs}$. In the separate coarse
limit $\Delta\downarrow0$, $\Delta/M^{1/5}\to\infty$,
\begin{equation}\label{eq:precision-coarse}
 \mathcal Q(M,\Delta)\sim
 K_{\rm coarse}M^{-1/4}\Delta^{1/4},\qquad
 K_{\rm coarse}=2^{-3/4}C_0\Bfun(1/2,1/8).
\end{equation}
Also $\mathcal H(\tau)\sim K_{\rm coarse}\tau^{1/4}$ as
$\tau\to\infty$. These limits do not depend on whether a fold lies
at a bin edge or inside a bin.
\end{theorem}

\subsection{A quadratic defect with an uncertain center}
For a finite nonnegative Radon measure $\nu$ on $\R$, define
\begin{align}
 \mathscr T_z(\nu)&=\sup_{f\in C_c^\infty(\R)}
   \left\{2\int_\R f\dd x-\int_\R |f'|^2\dd\nu
                         -\int_\R(x-z)^2f^2\dd x\right\},
                                          \label{eq:center-response}\\
 \mathcal F_{\rm ctr}(\eta)&=
   \inf_{\substack{d\ge0,\ d\in L^1(\R)\\\int_\R d=1}}
       \int_{-1/2}^{1/2}\mathscr T_{\eta v}(d\dd x)\dd v,
                        \qquad \eta\ge0. \label{eq:center-value}
\end{align}
Here $v$ in the second line is a scalar parameter. Thus $\eta>0$
means a center uniformly
distributed over an interval of width $\eta$, whereas $\eta=0$
means a known center. Measures will only be used for compactness;
the admissible mobility fields in \eqref{eq:center-value} are
densities. The response is an energy-dual quantity on an infinite
line, not a stationary transport coefficient there.

For a physical local coordinate with potential $a(x-z)^2$, allocated
mass $m$, and center uncertainty width $w$, set
$\ell=(m/a)^{1/5}$. The substitutions $x=\ell y$,
$D(x)=a\ell^4d(y)$, and $h(x)=(a\ell^2)^{-1}f(y)$ give the exact
conditional optimum
\begin{equation}\label{eq:center-scaling}
       a^{-4/5}m^{-1/5}
                    \mathcal F_{\rm ctr}\bigl(w/(m/a)^{1/5}\bigr).
\end{equation}
Translation and reflection preserve the mass constraint. Convexity
of the response permits an even $d$ in \eqref{eq:center-value}.

We first give the compactness fact needed to pass from this local
problem to the bin-dependent optimum. Its positive tail is a
proof device that can be added with arbitrarily small mass.

\begin{lemma}[Quadratic response under expanding-interval limits]
\label{lem:center-neumann}
Suppose $d\in L^1(\R)$ and
$d(x)\ge b(1+|x|)^{-\alpha}$ for some $b>0$ and $1<\alpha<2$.
Then $z\mapsto\mathscr T_z(d\dd x)$ is finite and continuous.
Let $L_n^-,L_n^+\to\infty$. On $(-L_n^-,L_n^+)$ take the
free-endpoint response with derivative energy $\int d|f'|^2$,
source $\int w_nf$, and reaction energy $\int w_n(x-z_n)^2f^2$.
If $z_n\to z$, the positive weights $w_n$ are uniformly bounded
above and below and converge locally uniformly to one, then the
interval response tends to $\mathscr T_z(d\dd x)$. The conclusion
is uniform over compact sets of center values, in the sequential
sense just stated, and is unchanged if the derivative coefficient
is multiplied by a scalar tending to one.
\end{lemma}
\begin{proof}
The proof parallels \cref{lem:graded-neumann-limit}, with a quadratic
tail in place of a quartic one; the slower source tail must be
retained. On a fixed core, a lower bound on $d$ and a subinterval
separated from all the possible centers control the $H^1$ norm by
the energy. Outside that core $(x-z_n)^2\ge c x^2$. Consequently,
writing $\mathfrak e_n$ for the interval energy,
\begin{equation}\label{eq:center-source-tail}
 \left|\int w_nf\right|\le C\mathfrak e_n[f]^{1/2},\qquad
 \int_{|x|>L}|w_nf|\dd x
                 \le CL^{-1/2}\mathfrak e_n[f]^{1/2}.
\end{equation}
The integrals are restricted to the interval where necessary.
Maximizing sequences can therefore be chosen with uniformly
bounded energy. They have locally weak $H^1$ limits and locally
weak weighted derivatives. Distributional differentiation
identifies the limiting derivative, and lower semicontinuity
controls its energy. The source integrals converge by
\eqref{eq:center-source-tail}.

To verify that the limit belongs to the whole-line smooth-test
completion, apply the one-dimensional Sobolev estimate on
$[x-1,x+1]$. Its local function and derivative squared norms are
bounded respectively by $C|x|^{-2}\mathfrak e[f]$ and
$C|x|^\alpha\mathfrak e[f]$. Thus
\begin{equation}\label{eq:center-point-bound}
 |f(x)|^2\le C\bigl(|x|^{-2}+|x|^{\alpha/2-1}\bigr)
                                \mathfrak e[f]
\end{equation}
for large $|x|$. In particular the limit is bounded. A cutoff
between $R$ and $2R$ adds derivative energy at most
$CR^{-2}\|f\|_\infty^2\int_{R<|x|<2R}d$, which vanishes; the
original energy tails vanish as well. On compact intervals,
approximate the derivative in $L^2(d\dd x)$ by smooth functions,
correct its integral with a fixed smooth bump when both traces
are prescribed, and integrate. The compact lower bound on $d$
makes this approximation uniform in the function, so source and
reaction also converge. This proves density in the required
energy space, including for unbounded $d$. The limiting value is
therefore no greater than the whole-line response. Conversely,
any compact smooth test is eventually admissible and its weighted
integrals converge, which proves the opposite inequality.

For continuity on the whole line, the same compact-core estimate
gives $\int(1+x^2)f^2\le C\mathfrak e_z[f]$ for centers in a
fixed compact set. Since
$|(x-z)^2-(x-z')^2|\le C|z-z'|(1+x^2)$ there, nearby forms
bound one another with factors tending to one. The source is
bounded in those energy norms by \eqref{eq:center-source-tail}.
Their response values are consequently continuous. The compactness
proof also applies along converging center sequences and to
scalar changes of the derivative coefficient, proving the final
assertions.
\end{proof}

\begin{proposition}[Local value and its endpoint limits]
\label{prop:center-value}
For every $\eta\ge0$, the infimum in \eqref{eq:center-value} is
attained by a nonnegative $L^1$ density of mass one; an even
minimizer can be chosen. The value is continuous on $[0,\infty)$
and satisfies
\begin{equation}\label{eq:center-bounds}
 \max\{C_{\rm pl},C_0\eta^{1/4}\}
 \le\mathcal F_{\rm ctr}(\eta)\le C(1+\eta^{1/4}),
 \qquad
 \mathcal F_{\rm ctr}(0)=C_{\rm pl}.
\end{equation}
Moreover,
\begin{equation}\label{eq:center-large}
       \mathcal F_{\rm ctr}(\eta)\sim C_0\eta^{1/4}
                         \quad(\eta\to\infty).
\end{equation}
No uniqueness or monotonicity in the interval width is needed
for these assertions or for \cref{thm:precision}.
\end{proposition}
\begin{proof}
Allowing the center to be known can only reduce the objective.
The translated quadratic placement theorem therefore gives
$\mathscr T_z(d\dd x)\ge C_{\rm pl}$ for every unit-mass $d$
and every $z$, including an infinite response. This proves the
first lower bound and equality at zero.

For the other lower bound fix $\psi\in C_c^\infty(\R)$ and write
$I_\psi=\int\psi$, $T_\psi=\int|\psi'|^2$,
$U_\psi=\int x^2\psi^2$. For $\eta>0$ use the translated tests
\[
 h_z(x)=\eta^{1/2}\psi((x-z)/\eta^{-1/4}),
                       \qquad |z|\le\eta/2.
\]
Their source minus reaction is
$\eta^{1/4}(2I_\psi-U_\psi)$, independently of $z$.
At every $x$, the averaged squared derivative is bounded by
\[
 \frac1\eta\int_{-\eta/2}^{\eta/2}|h_z'(x)|^2\dd z
       \le\eta^{1/4}T_\psi.
\]
Integrating this bound against any unit-mass mobility gives
$\mathcal F_{\rm ctr}(\eta)\ge
\eta^{1/4}(2I_\psi-U_\psi-T_\psi)$. Take the oscillator
variational supremum to obtain $C_0\eta^{1/4}$. This lower
certificate allows arbitrary mobility spikes and zero sets.

For an upper bound put constant mobility $1/(\eta+2)$ on
$[-\eta/2-1,\eta/2+1]$ and zero outside. Neumann bracketing
and reciprocal-potential integration outside this interval give
a bound $C(1+\eta^{1/4})$. When $\eta\to\infty$, the harmonic
length $(\eta+2)^{-1/4}$ is much smaller than the unit padding.
The natural-endpoint oscillator limit is uniform over all center
positions, so the interior contribution is
$C_0(\eta+2)^{1/4}[1+o(1)]$. The exterior contribution is at
most two. This proves \eqref{eq:center-large}.

For compactness, the singular-mass removal proof in
\cref{lem:measure-relaxation} applies to \eqref{eq:center-response}
without change: derivative flattening near a null support changes
the source and the locally bounded reaction by $o(1)$. Hence
\begin{equation}\label{eq:center-singular}
      \mathscr T_z(d\dd x+\nu_s)=\mathscr T_z(d\dd x).
\end{equation}
For a compact smooth test its value is continuous under vague
convergence of finite measures and simultaneous convergence of
$z$. A common countable set of compact smooth tests, on a
countable exhaustion of the line, gives measurability of the
supremum. That supremum is jointly lower semicontinuous.

Take a minimizing sequence of unit-mass densities for fixed
$\eta$. It has a vaguely converging subsequence with limit $\nu$
of mass at most one. Fatou's lemma, after the fixed parameterization
$z=\eta v$, bounds the objective at $\nu$ by the liminf. Remove
its singular part using \eqref{eq:center-singular}, and add a
nonnegative density to replace any missing mass. Increasing the
derivative coefficient cannot increase the response. The resulting
unit-mass density therefore has value no greater than the infimum
and is a minimizer. This argument makes no assumption that the
original minimizing sequence was tight, or that its weak limit
was absolutely continuous. Convex reflection averaging gives an
even minimizer. The same argument with $\eta_n\to\eta$ proves
lower semicontinuity of $\mathcal F_{\rm ctr}$.

For the upper semicontinuity, choose a unit-mass density $d$ with
objective within $\varepsilon$ of the value at $\eta$. Let
$\rho_\alpha$ be a positive unit-mass multiple of
$(1+|x|)^{-\alpha}$, $1<\alpha<2$, and set
\begin{equation}\label{eq:center-regularization}
       d_\varepsilon=\frac{d+\varepsilon\rho_\alpha}{1+\varepsilon}.
\end{equation}
It has mass one and a positive integrable tail. Form ordering
gives $\mathscr T_z(d_\varepsilon\dd x)\le
(1+\varepsilon)\mathscr T_z(d\dd x)$.
By \cref{lem:center-neumann}, the response of this fixed
regularization is continuous, uniformly for centers in a compact
set. Its average is thus continuous in $\eta$. Use it as a
competitor at $\eta_n$, and then let $\varepsilon\downarrow0$.
This proves upper semicontinuity and completes the proof.
\end{proof}

\subsection{Uniform localization of a conditional optimum}
For a bin $I$ with midpoint $c_I\in(-1,1)$, set
\begin{align}
 a_I&=1-c_I^2,\quad r_I=\arccos(-c_I),\quad m=M/2,
                \quad\ell_I=(m/a_I)^{1/5},\notag\\
 \eta_I&=\frac{\Delta}{\sqrt{a_I}\ell_I}
          =2^{1/5}\frac{\Delta}{M^{1/5}}a_I^{-3/10}.
                         \label{eq:bin-local-scales}
\end{align}
Let $V(M,I)$ be the infimum of the conditional mean
$\Delta^{-1}\int_IJ_c(D)\dd c$ over all budget-$M$ fields.

\begin{lemma}[Conditional variational localization]
\label{lem:conditional-localization}
On every fixed compact $K\subset(-1,1)$, uniformly for bins with
$c_I\in K$ as $M,\Delta\downarrow0$ and
$\Delta/M^{1/5}$ remains bounded,
\begin{equation}\label{eq:conditional-limit}
 \frac{V(M,I)}{2a_I^{-4/5}(M/2)^{-1/5}}
                    -\mathcal F_{\rm ctr}(\eta_I)\longrightarrow0.
\end{equation}
The lower bound allows every nonnegative $L^1$ field. The upper
bound is realized by a field with exactly mass $M$ selected using
only the bin, and can be made uniform on $K$.
\end{lemma}
\begin{proof}
Reflection $s\mapsto-s$ preserves $k_c$ for each $c$, not just
after ensemble averaging. Convexity therefore allows replacement
of any competitor by $[D(s)+D(-s)]/2$ without increasing its
conditional cost. The two half-circles then have exactly mass
$m=M/2$. This symmetry does not require that $I$ be symmetric
under $c\mapsto-c$.

For the lower bound consider any sequence of bins as in the
statement, and choose a subsequence attaining the liminf of the
normalized conditional values. By further subsequences,
$c_I\to c_*\in K$ and $\eta_I\to\eta_*$. Take symmetric
competitors whose normalized costs approach that liminf. Push
their mass on $(0,\pi)$ to $x=(s-r_I)/\ell_I$, dividing it by
$m$. These probability measures on expanding intervals have a
vaguely converging subsequence with limit $\nu$, of mass at most
one. For $c=c_I+\Delta v$, $|v|\le1/2$, the scaled potential
satisfies, locally uniformly in $x,v$,
\[
 \frac{c+\cos(r_I+\ell_I x)}{\sqrt{a_I}\ell_I}
              =\eta_Iv-x+O_K(\ell_Ix^2).
\]
Insert any compact smooth test at $r_I$, with amplitude
$(a_I\ell_I^2)^{-1}$, and its reflected copy at $-r_I$.
The supports are disjoint for small $M$. The scaled source and
reaction converge to those in \eqref{eq:center-response}, and
the derivative integral converges vaguely. Taking the supremum
over compact tests gives, for each $v$,
\begin{equation}\label{eq:conditional-liminf}
 \liminf\frac{J_{c_I+\Delta v}(D_I)}{2a_I^{-1}\ell_I^{-1}}
                         \ge\mathscr T_{\eta_*v}(\nu).
\end{equation}
Fatou's lemma then gives the averaged lower bound
$\int_{-1/2}^{1/2}\mathscr T_{\eta_*v}(\nu)\dd v
\ge\mathcal F_{\rm ctr}(\eta_*)$: remove singular mass and
complete the remaining density to mass one as in
\cref{prop:center-value}. Continuity of the local value turns
this sequential bound into the uniform lower assertion. No
regularity, scale separation within the design, or absence of
escaped mass has been imposed on the competitor.

For recovery fix a limit pair $c_*,\eta_*$ and a unit-mass
profile $d_\varepsilon$ as in \eqref{eq:center-regularization},
whose averaged value at $\eta_*$ is within an arbitrarily small
error of $\mathcal F_{\rm ctr}(\eta_*)$. In a fixed small
neighborhood of $r_I$, use the exact coordinate
\begin{equation}\label{eq:regular-exact-coordinate}
 x=\frac{-\cos s-c_I}{\sqrt{a_I}\ell_I},\qquad
 \dd s=\ell_Iw_I(x)\dd x,\qquad
 w_I(x)=\frac{\sqrt{a_I}}{\sin s}.
\end{equation}
The neighborhood stays inside $(0,\pi)$ uniformly on $K$.
Its scaled endpoints tend to both infinities; $w_I$ is uniformly
bounded above and below and tends locally uniformly to one.
For every $c$ in the bin,
\[
 k_c(s)=a_I\ell_I^2(x-\eta_Iv)^2,
                         \qquad c=c_I+\Delta v,
\]
exactly. Put the preliminary field
\begin{equation}\label{eq:bin-recovery}
        D_I(s)=a_I\ell_I^4w_I(x)d_\varepsilon(x)
\end{equation}
on this neighborhood, its reflection on the other root
neighborhood, and zero elsewhere. Its mass on each neighborhood
is $m\int w_I^2d_\varepsilon\dd x=m[1+o(1)]$ by dominated
convergence. A single scalar normalization therefore makes its
total mass exactly $M$ and tends to one.

Under the response scaling $(a_I\ell_I)^{-1}$, the derivative
energy in \eqref{eq:bin-recovery} is precisely
$\int d_\varepsilon|f'|^2$; source and reaction have the weight
$w_I$. Apply Neumann bracketing at the neighborhood endpoints
and \cref{lem:center-neumann}. The exterior reciprocal-potential
integral is bounded on $K$, because all possible roots lie
strictly inside the two neighborhoods. Its contribution vanishes
after multiplication by $a_I\ell_I$. Consequently this recovery
has normalized conditional value tending to
$\int_{-1/2}^{1/2}\mathscr T_{\eta_*v}(d_\varepsilon\dd x)\dd v$.
Normalization changes the response by a factor tending to one,
as the full forms are bounded by the minimum and maximum of one
and that normalization factor. Let $\varepsilon\downarrow0$.

These arguments prove the upper assertion along every converging
sequence in the compact parameter set. Equivalently, one may
choose a finite collection of regularized profiles: continuity
of their averaged responses supplies a finite cover of the
compact range of $\eta_I$. Assign the first qualifying profile
to each bin and use \eqref{eq:bin-recovery}. This gives uniform
errors, a deterministic bin policy, and an exact per-bin budget.
Together with the sequential lower bound it proves
\eqref{eq:conditional-limit}.
\end{proof}

\subsection{Global bounds, including bins crossing a fold}
The regular-bin limit is insufficient by itself: curvatures vanish
at $c=\pm1$, and some bins contain both rootless and two-root
realizations. We record a single family of admissible trials that
controls all these bins uniformly.

\begin{lemma}[Bin policies with uniform integrable bounds]
\label{lem:bin-envelope}
Put $W=\Delta+M^{2/7}$. There are fixed constants such that bins
on the two-root side at distance $t=1-|c|\ge C W$ from the
nearest fold admit a common-within-bin field satisfying
\begin{equation}\label{eq:bin-regular-envelope}
 J_c(D_I)\le C\left[M^{-1/5}t^{-4/5}
                    +M^{-1/4}\Delta^{1/4}t^{-7/8}\right]
                         \quad(c\in I).
\end{equation}
All bins intersecting a fixed multiple of the fold neighborhoods
$|1-|c||\le C W$ can be covered by two groups. Each group admits
a single budget-$M$ field whose integrated response over that
group is at most
\begin{equation}\label{eq:bin-fold-envelope}
                  C M^{-1/4}W^{3/8}.
\end{equation}
The integrated cost of the remaining rootless bins is at most
$C W^{-1/2}$ and is absorbed in \eqref{eq:bin-fold-envelope}.
\end{lemma}
\begin{proof}
In a regular bin, $t$ is comparable to its midpoint value.
Each possible-root arc has length at most $Cw$, where
$w=\Delta/\sqrt t$. Pad each arc by a fixed multiple of
$\ell=(M/t)^{1/5}$ and put half the budget uniformly on each
padded arc. The patches have length comparable to $w+\ell$,
and their mobility $e$ is comparable to $M/(w+\ell)$. Since
$t\ge C W$, increasing the fixed $C$ if necessary makes the
patches disjoint and keeps them in quadratic neighborhoods:
\[
       \frac w{\sqrt t}=\frac\Delta t\ll1,\qquad
       \frac\ell{\sqrt t}=(M/t^{7/2})^{1/5}\ll1.
\]
The constants denoted by $\ll1$ here are fixed sufficiently
small constants, uniformly over the construction. The harmonic
length $(e/t)^{1/4}$ is at most a constant times $\ell$, so every
possible root has a fixed positive margin in harmonic units.
The interval oscillator bound and quadratic potential comparison
give an interior response $Ce^{-1/4}t^{-3/4}$. Outside the
patches, reciprocal-potential integration gives
$C/(t\ell)+Ct^{-3/2}$, which is bounded by the same scale.
Expanding $(w+\ell)^{1/4}$ proves
\eqref{eq:bin-regular-envelope}.

Classify as fold bins every bin meeting the stated fold
neighborhood. Since a bin has width $\Delta\le W$, their
union remains inside a possibly larger fixed multiple of that
neighborhood, irrespective of grid alignment. Around the
corresponding fold site $s_j$ put constant mobility on an interval
of radius $R=K_0\sqrt W$, with fixed $K_0$ large enough to
contain every possible root with a fixed relative margin.
Normalize it to mass $M$, so $e\asymp M/\sqrt W$, and put
zero mobility outside. In the exact sine coordinate
$z=2\sin((s-s_j)/2)$ used in the proof of
\cref{lem:cosine-uniform}, the local potential is
$(z^2/2-t)^2$, with uniformly comparable metric factors.
The inequality $W\ge M^{2/7}$ ensures $W\ge c e^{1/3}$.
The quartic inner layer and both tails from
\cref{lem:cosine-uniform} therefore give
\begin{align*}
 \int_{|t|\le C W}J_{\rm local}(t)\dd t
 &\le C\left[e^{-1/6}
           +e^{-1/4}\int_{e^{1/3}}^{C W}t^{-3/4}\dd t\right]\\
 &\le C e^{-1/4}W^{1/4}
       \le C M^{-1/4}W^{3/8}.
\end{align*}
An empty integral is omitted, and fixed comparison factors in
the lower integration limit are harmless. The same bounds hold
for the finite interval with free endpoints: each scaled root
is separated from its endpoint by a fixed fraction of the
interval radius. They follow from the interval bracketing and
localization proof of \cref{lem:cosine-uniform}, not from an
unrestricted large-parameter whole-line limit.

Outside the patch the reciprocal-potential response is at most
$CW^{-3/2}$ per offset, hence $CW^{-1/2}$ after integration.
This is absorbed above because
$M^{1/4}W^{-7/8}\le1$. Finally, on the rootless side beyond
these bins, $J_c(D)\le\int_\Gamma k_c^{-1}$ for any field.
The rootless reciprocal-potential bound integrates to
$CW^{-1/2}+C$, also absorbed. All trial fields have exactly
the prescribed mass, including those reused over an entire
group of bins.
\end{proof}

We can now prove the order theorem without imposing any scale
condition on the relative precision. The lower bound
$\mathcal Q(M,\Delta)\ge\mathcal O(M)\ge cM^{-1/5}$ follows
from exact observation. For the coarse contribution retain bins
in $[-1/2,1/2]$ and assume $\Delta\ge L M^{1/5}$ with $L$
a sufficiently large fixed number. In each retained bin the
possible roots range over arcs of length comparable to $\Delta$,
with conditional densities at most $C/\Delta$. Use compact
nonnegative translated bumps at these roots, of width
$b=(M/\Delta)^{1/4}$ and amplitude $\gamma b^{-2}$.
Their supports lie in uniformly quadratic neighborhoods, and
the two root supports are disjoint. Averaging source minus
reaction gives at least $c\gamma b^{-1}-C\gamma^2b^{-1}$.
At any spatial point the averaged squared derivative is at
most $C\gamma^2/(\Delta b^5)$. Against an arbitrary total
mass $M$ its cost is at most $C\gamma^2b^{-1}$. Choose a
fixed sufficiently small $\gamma>0$. This proves a conditional
lower bound $cM^{-1/4}\Delta^{1/4}$ for every field in every
retained bin. Their total probability is bounded below as
$\Delta\to0$. When $\Delta<L M^{1/5}$ this term is already
controlled by the oracle lower bound. The maximum of the two
bounds is comparable to their sum.

Integrating \eqref{eq:bin-regular-envelope} gives the two terms
in \eqref{eq:precision-order}, since $4/5$ and $7/8$ are both
less than one. The remaining terms satisfy
\begin{equation}\label{eq:fold-smallness}
 M^{-1/4}W^{3/8}
       \le C\bigl(M^{-1/7}+M^{-1/4}\Delta^{3/8}\bigr)
       =o\bigl(M^{-1/5}+M^{-1/4}\Delta^{1/4}\bigr).
\end{equation}
This proves the joint order law. For any fixed $N$ the separate
statement is $\mathcal Q(M,4/N)\asymp M^{-1/4}$, with constants
allowed to depend on $N$. A positive-length regular portion of
one bin gives the same moving-root lower bound, and the
predetermined mean design from \cref{thm:blind-sharp} supplies
the upper bound. No small-$\Delta$ coefficient is asserted for
that fixed partition.

\subsection{Assembly of the exact crossover and sharp coarse limit}
Suppose first that $\Delta/M^{1/5}\to\tau<\infty$. On a fixed
compact subinterval of $(-1,1)$, \cref{lem:conditional-localization}
and the finite-bin decomposition of \eqref{eq:bin-value} give
Riemann sums with density
\begin{equation}\label{eq:crossover-density}
 \frac{2^{6/5}}4(1-c^2)^{-4/5}
 \mathcal F_{\rm ctr}\!\left(2^{1/5}\tau(1-c^2)^{-3/10}\right).
\end{equation}
Uniformity in that lemma, and continuity of the local value,
justify these sums even though the designs and number of bins
both change. Discarding all other bins gives the global liminf;
then enlarge the compact interval to $(-1,1)$.

For the matching upper bound, use the recovery policies on the
retained bins and the policies of \cref{lem:bin-envelope}
elsewhere. Any bins cut by the artificial ends of the retained
interval have width at most $\Delta$ and can be assigned to the
omitted set. The normalized regular-bin envelope is
\[
        Ct^{-4/5}
           +C(\Delta/M^{1/5})^{1/4}t^{-7/8},
\]
whose integral over $0<t<\delta$ is at most
$C\delta^{1/5}+C\sup(\Delta/M^{1/5})^{1/4}\delta^{1/8}$.
The normalized fold contribution tends to zero. Indeed
$M^{1/5}M^{-1/7}=M^{2/35}\to0$ and, for bounded
$\Delta/M^{1/5}$,
$M^{-1/20}\Delta^{3/8}=O(M^{1/40})\to0$.
Send $M,\Delta\to0$ first and then $\delta\downarrow0$ to
prove \eqref{eq:precision-crossover}, including $\tau=0$.
The bound \eqref{eq:center-bounds} also shows that the integrand
in \eqref{eq:crossover-function} is bounded on bounded $\tau$
sets by an integrable multiple of
$(1-c^2)^{-4/5}+(1-c^2)^{-7/8}$. Dominated convergence proves
continuity of $\mathcal H$ and
\[
 \mathcal H(0)=\frac{2^{6/5}C_{\rm pl}}4\Bfun(1/2,1/5)
                    =K_{\rm obs}.
\]

The coarse limit requires its own uniform argument; letting
$\tau\to\infty$ after a fixed-$\tau$ theorem would not prove
a simultaneous limit. Retain again a compact regular offset
set. In a bin with midpoint $c_I$, each root arc has length
$w=\Delta/\sqrt{a_I}\,[1+o(1)]$, uniformly on this set.
Put $e=M/(2w)$ and $b=(e/a_I)^{1/4}$. The coarse hypothesis
is equivalent here to $b/w\to0$. Take harmonic tests centered
at the two actual roots, with amplitude $(e a_I)^{-1/2}$,
scale $b$, and a fixed compact smooth shape $\psi$. Set
$z_0=e^{-1/4}a_I^{-3/4}$. Their averaged source minus
reaction is
\[
       2z_0(2I_\psi-U_\psi)+o(z_0).
\]
The conditional root density on each arc is
$(1+o(1))/w$. The two neighborhoods are disjoint, so at every
spatial point the sum of the two averaged derivative kernels
is at most
\[
       [1+o(1)]\frac{T_\psi}{w}e^{-5/4}a_I^{-3/4}.
\]
Bin edges only truncate this nonnegative kernel. Multiplication
by any competing mass $M=2ew$ gives a derivative cost at most
$[1+o(1)]2z_0T_\psi$. First take the joint limit with $\psi$
fixed, then take its oscillator variational supremum. This proves
the unrestricted conditional lower bound $[1-o(1)]2C_0z_0$.
This proof does not presume an allocation by the competitor.

For the upper bound, divide the mass equally over the two root
arcs, with padding $p$ chosen so that $b\ll p\ll w$; for
example $p=\sqrt{bw}$ works uniformly on the retained set.
The padded constant mobility is $e[1+o(1)]$. Neumann harmonic
localization is uniform because every root is many harmonic
lengths from each endpoint. The exterior reciprocal-potential
contribution is $O((a_Ip)^{-1})=o(z_0)$. Curvature and the
offset-to-root Jacobian vary by $o(1)$ on these shrinking arcs.
Both bounds therefore give
\begin{equation}\label{eq:coarse-conditional}
 V(M,I)=[1+o(1)]2^{5/4}C_0a_I^{-7/8}
                               M^{-1/4}\Delta^{1/4}
\end{equation}
uniformly on every fixed regular offset set.

After dividing \eqref{eq:bin-regular-envelope} by the coarse
scale, the omitted regular-bin integrand is bounded by
$C(M^{1/5}/\Delta)^{1/4}t^{-4/5}+Ct^{-7/8}$.
In this regime $W\sim\Delta$, so the normalized fold term is
$O(\Delta^{1/8})\to0$. Riemann sums in
\eqref{eq:coarse-conditional} thus extend to the full ensemble,
with coefficient
\[
       \frac{2^{5/4}C_0}4\int_{-1}^1(1-c^2)^{-7/8}\dd c
          =2^{-3/4}C_0\Bfun(1/2,1/8).
\]
This proves \eqref{eq:precision-coarse}. Finally, applying
\eqref{eq:center-large} in \eqref{eq:crossover-function} and
using the same integrable bound gives
$\mathcal H(\tau)/\tau^{1/4}\to K_{\rm coarse}$.
This last consistency statement follows from the local value;
the simultaneous coarse limit has been proved separately.

\subsection{Transport consequence and meaning of the resolution scale}
For the full-bulk optimum $\mathcal Q_{\rm bulk}$ with the same
bin observation, \cref{cor:all-information} gives, in the
dimensionless scalar convention of this paper,
\begin{equation}\label{eq:precision-bulk}
 1\le\frac{\mathcal Q_{\rm bulk}(M,\Delta)}
                {\chi\mathcal Q(M,\Delta)}
       \le1+CM^{1/5}\log(1/M)
\end{equation}
for small $M$, uniformly over the partitions under consideration.
The assumptions are those of \cref{thm:optimal-transfer}, including
fixed positive bulk transverse diffusion and nonzero mean speed.
Mixing each selected field with the same vanishing uniform
background preserves both its bin dependence and its exact
budget. Thus every scalar equivalent and order above transfers
to optimized flow-induced dispersion. In physical units the
response factor $\ell_{\rm ref}/k_{\rm ref}$ in
\eqref{eq:J-dimensional} must also be included. This is an
optimal-value comparison, not a remainder estimate for every
unmixed trial.

For one regular defect the uncertain root span is
$w\simeq\Delta/\sqrt a$. Equation \eqref{eq:center-scaling}
places the local crossover at
\begin{equation}\label{eq:local-precision-threshold}
             \Delta\asymp m^{1/5}a^{3/10}.
\end{equation}
The global mean is dominated by ordinary roots, whose curvatures
give an integrable coefficient, so its threshold is
$\Delta\asymp M^{1/5}$. For $N=2^B$ equal bins, this means
$B=(1/5)\log_2(1/M)+O(1)$ binary digits of offset resolution.
In the sharp coarse regime, one further binary digit multiplies
the leading response by $2^{-1/4}$. Refining a nested partition
cannot worsen its optimum, since every coarser policy remains
available; arbitrary nonnested partitions are not ordered by
their widths alone.

These statements concern noiseless quantization followed by a
static mobility choice. They do not optimize a measurement
channel, assign a cost to information, or impose a fabrication
length or mobility cap. The local value supplies an attained
variational coefficient for the intermediate precision regime;
the theorem does not require an explicit formula for its
minimizing profile.
